\let\negmedspace\undefined
\let\negthickspace\undefined
\documentclass[journal]{IEEEtran}
\usepackage[a5paper, margin=10mm, onecolumn]{geometry}
%\usepackage{lmodern} % Ensure lmodern is loaded for pdflatex
\usepackage{tfrupee} % Include tfrupee package

\setlength{\headheight}{1cm} % Set the height of the header box
\setlength{\headsep}{0mm}     % Set the distance between the header box and the top of the text

\usepackage{gvv-book}
\usepackage{gvv}
\usepackage{cite}
\usepackage{amsmath,amssymb,amsfonts,amsthm}
\usepackage{algorithmic}
\usepackage{graphicx}
\usepackage{textcomp}
\usepackage{xcolor}
\usepackage{txfonts}
\usepackage{listings}
\usepackage{enumitem}
\usepackage{mathtools}
\usepackage{gensymb}
\usepackage{comment}
\usepackage[breaklinks=true]{hyperref}
\usepackage{tkz-euclide} 
\usepackage{listings}
% \usepackage{gvv}                                        
\def\inputGnumericTable{}                                 
\usepackage[latin1]{inputenc}                                
\usepackage{color}                                            
\usepackage{array}                                            
\usepackage{longtable}                                       
\usepackage{calc}                                             
\usepackage{multirow}                                         
\usepackage{hhline}                                           
\usepackage{ifthen}                                           
\usepackage{lscape}
\begin{document}
\bibliographystyle{IEEEtran}
\title{12.744}
\author{EE25BTECH11002 - Achat Parth Kalpesh }
{\let\newpage\relax\maketitle}
\renewcommand{\thefigure}{\theenumi}
\renewcommand{\thetable}{\theenumi}
\setlength{\intextsep}{10pt} % Space between text and floats
\numberwithin{equation}{enumi}
\numberwithin{figure}{enumi}
\renewcommand{\thetable}{\theenumi}
\parindent 0px



\textbf{Question:}\\
Consider the following equations of straight lines:
\begin{align*}
	\text{Line L1} &\colon 2x - 3y = 5 \\
	\text{Line L2} &\colon 3x + 2y = 8 \\
	\text{Line L3} &\colon 4x - 6y = 5 \\
	\text{Line L4} &\colon 6x - 9y = 6
\end{align*}
Which one among the following is the correct statement?
\begin{enumerate}
	\item L1 is parallel to L2 and L1 is perpendicular to L3
	\item L2 is parallel to L4 and L2 is perpendicular to L1
	\item L3 is perpendicular to L4 and L3 is parallel to L2
	\item L4 is perpendicular to L2 and L4 is parallel to L3
\end{enumerate}

\textbf{Solution:}\\
A line $ax + by = c$ can be written as
\begin{align}
\vec{n}^T \vec{x} = c,
\end{align}
where
\begin{align}
\vec{n} = \myvec{a\\b} \quad
\vec{x} = \myvec{x\\y}
\end{align}

Thus, the normal vectors are:
\begin{align}
\vec{n}_1 &= \myvec{2\\-3}, &
\vec{n}_2 &= \myvec{3\\2}, &
\vec{n}_3 &= \myvec{4\\-6}, &
\vec{n}_4 &= \myvec{6\\-9}
\end{align}

Two lines are parallel if their normal vectors are proportional:
\begin{align}
\vec{n}_3 &= 2 \vec{n}_1 \implies \text{L1 is parallel to L3}, \\
\vec{n}_4 &= \frac{3}{2} \vec{n}_3 \implies \text{L3 is parallel to L4}.
\end{align}

Two lines are perpendicular if
\begin{align}
\vec{n}_i^T \vec{n}_j = 0
\end{align}

\begin{align}
\vec{n}_1^T \vec{n}_2 &= \myvec{2 & -3} \myvec{3\\2} = 6 - 6 = 0 
\implies \text{L1 is perpendicular to L2},\\
\vec{n}_4^T \vec{n}_2 &= \myvec{6 & -9} \myvec{3\\2} = 18 - 18 = 0 
\implies \text{L4 is perpendicular to L2}
\end{align}

L4 is perpendicular to L2 and L4 is parallel to L3

\begin{figure}[h!]
    \centering
    \includegraphics[width=\columnwidth]{figs/figure_py.png}
    \caption{Lines}
    \label{fig:fig}
 \end{figure}
\end{document}